\begin{frame}{ADC --- Conversor Analógico para Digital}
    ADC significa \textit{Analog-to-Digital Converter}.

    \vspace{0.2cm}

    Sua função é transformar uma tensão analógica em um número binário.

    \begin{center}
        \resizebox{0.9\linewidth}{!}{%
        \begin{tikzpicture}[
            bloco/.style={draw, rounded corners, minimum width=3cm,
                          minimum height=1.2cm, align=center, fill=blue!10},
            seta/.style={-{Latex[length=3mm]}, thick}
        ]
            \node[bloco] (entrada) {Entrada analógica\\$V_{in}$};
            \node[bloco, right=1.5cm of entrada] (adc) {ADC};
            \node[bloco, right=1.5cm of adc] (saida) {Saída digital\\$b_{n-1}\ldots b_1b_0$};

            \draw[seta] (entrada) -- (adc);
            \draw[seta] (adc) -- (saida);
        \end{tikzpicture}%
        }
    \end{center}

    \begin{exampleblock}{Exemplo}
        Um sensor de temperatura fornece $1{,}65$ V. O ADC mede essa tensão
        e a representa como um valor numérico para o microcontrolador.
    \end{exampleblock}
\end{frame}